% ============================================================
% section1_preparation.tex - Раздел 1: Подготовка инструментов
% ============================================================
% !TeX program = xelatex
% Компиляция: xelatex --shell-escape main.tex

\section{Подготовка инструментов}

Добро пожаловать в мир программирования для инженеров! Прежде чем мы начнём обрабатывать терагерцовые спектры, нам нужно подготовить рабочее место. В этом разделе мы установим язык Python, выберем среду разработки (IDE) и настроим проект. Не волнуйся, если ты никогда этого не делал — мы пройдём каждый шаг вместе.

\subsection{Установка Python}
\subsubsection{Скачивание и установка}

Python — это язык программирования, который будет выполнять наши команды. Чтобы компьютер начал его понимать, нужно установить интерпретатор.

Перейди на официальный сайт \url{https://www.python.org/downloads/} и скачай последнюю версию для Windows. Запусти скачанный установщик.

\warning{При установке обязательно поставь галочку \textbf{«Add Python to PATH»} (Добавить Python в PATH) на первом экране установщика. Это позволит запускать скрипты из любой папки в командной строке.}

\smartfigure{fig1.1.png}{Окно установщика Python с выделенной галочкой «Add Python to PATH».}{fig:1_1}

\subsubsection{Проверка установки}

Давай убедимся, что всё установилось правильно. Открой командную строку (нажми \texttt{Win + R}, введи \texttt{cmd} и нажми Enter).

Введи следующую команду и нажми Enter:

\begin{minipage}{\linewidth}
	\captionof{listing}{Проверка версии Python}
	\begin{minted}{bash}
python --version
	\end{minted}
\end{minipage}

Если ты видишь ответ вроде \texttt{Python 3.12.x}, значит всё отлично!

\subsection{Выбор и установка среды разработки}
\subsubsection{Обзор вариантов}

Писать код можно даже в обычном Блокноте, но это неудобно. Специализированные программы — среды разработки (IDE) — подсвечивают ошибки, подсказывают команды и сильно упрощают жизнь. 

Для нашего проекта мы выберем \textbf{PyCharm Community Edition}. Это бесплатная, мощная и очень популярная среда от компании JetBrains. Она отлично подходит как для новичков, так и для профессионалов.

\subsubsection{Установка PyCharm}

Скачай установщик со страницы \url{https://www.jetbrains.com/pycharm/download/} (выбирай версию Community). Установка стандартная: просто нажимай «Далее».

\smartfigure{fig1.2.png}{Стартовое окно PyCharm после первого запуска.}{fig:1_2}

\subsection{Виртуальное окружение}
\subsubsection{Зачем оно нужно}

Представь, что ты работаешь над двумя проектами: один требует старую версию библиотеки для графиков, а другой — новую. Если установить их в систему глобально, они могут конфликтовать. 

\textbf{Виртуальное окружение (venv)} — это изолированная папка для конкретного проекта. В ней хранится свой Python и свои библиотеки. Это правило хорошего тона в мире программирования: один проект = одно виртуальное окружение.

\subsubsection{Создание и активация}

В PyCharm создание виртуального окружения происходит почти автоматически при создании нового проекта (мы сделаем это в следующем разделе). Но полезно знать, как это делается в командной строке:

\begin{minipage}{\linewidth}
	\captionof{listing}{Создание и активация виртуального окружения}
	\begin{minted}{bash}
# Создание папки .venv с виртуальным окружением
python -m venv .venv

# Активация окружения
.venv\Scripts\activate
	\end{minted}
\end{minipage}

\note{Когда окружение активировано, в начале строки терминала появится приписка \texttt{(.venv)}.}

\subsection{Установка библиотек}
Голый Python умеет многое, но для научных расчётов нам понадобятся дополнительные инструменты (библиотеки). Мы будем использовать:
\begin{itemize}
    \item \textbf{NumPy} — для быстрых математических операций и работы с массивами.
    \item \textbf{SciPy} — для сложных вычислений (БПФ, фильтрация, оптимизация).
    \item \textbf{Matplotlib} — для построения красивых графиков.
\end{itemize}

\subsubsection{Создание файла requirements.txt}

Вместо того чтобы устанавливать библиотеки по одной, создают файл \texttt{requirements.txt}, где перечислены все зависимости. Это удобно, если ты захочешь запустить код на другом компьютере.

\begin{minipage}{\linewidth}
	\captionof{listing}{Содержимое файла requirements.txt}
	\begin{minted}{text}
numpy>=1.24.0
scipy>=1.10.0
matplotlib>=3.7.0
	\end{minted}
\end{minipage}

\subsubsection{Установка через pip}

\texttt{pip} — это менеджер пакетов Python. Чтобы установить всё из нашего списка, выполни команду в терминале (убедись, что виртуальное окружение активировано!):

\begin{minipage}{\linewidth}
	\captionof{listing}{Команда установки библиотек}
	\begin{minted}{bash}
pip install -r requirements.txt
	\end{minted}
\end{minipage}

Компьютер скачает и установит нужные файлы.

\subsection{Проверка работоспособности}

Давай напишем наш первый скрипт, чтобы убедиться, что всё работает как часы. Создай файл \texttt{check.py} и скопируй в него этот код:

\begin{minipage}{\linewidth}
	\captionof{listing}{Скрипт для проверки библиотек (check.py)}
	\begin{minted}{python}
import sys
import numpy as np
import scipy
import matplotlib

# Выводим версии установленных компонентов
print(f"Версия Python: {sys.version.split()[0]}")
print(f"Версия NumPy: {np.__version__}")
print(f"Версия SciPy: {scipy.__version__}")
print(f"Версия Matplotlib: {matplotlib.__version__}")

print("\nОтлично! Все инструменты готовы к работе.")
	\end{minted}
\end{minipage}

\tip{Запусти этот скрипт. Если в консоли появились версии библиотек и ободряющее сообщение — поздравляю, ты справился с настройкой!}

\subsection{Итоги раздела}
Мы проделали большую работу:
\begin{enumerate}
    \item Установили Python и добавили его в системный путь (PATH).
    \item Выбрали среду разработки PyCharm.
    \item Узнали, зачем нужны виртуальные окружения.
    \item Установили джентльменский набор инженера: NumPy, SciPy и Matplotlib.
\end{enumerate}

В следующем разделе мы создадим структуру нашего проекта и заложим фундамент для будущей программы. Переворачивай страницу, когда будешь готов!